\documentclass{article}

\usepackage{amsmath}
\usepackage{listings}
\usepackage{color}
\usepackage{hyperref}
\usepackage{geometry}
\geometry{margin=1in}

% Define code styling
\definecolor{codegray}{rgb}{0.5,0.5,0.5}
\definecolor{codepurple}{rgb}{0.58,0,0.82}
\definecolor{backcolour}{rgb}{0.95,0.95,0.92}

\lstdefinestyle{mystyle}{
    backgroundcolor=\color{backcolour},   
    commentstyle=\color{codegray},
    keywordstyle=\color{blue},
    numberstyle=\tiny\color{codegray},
    stringstyle=\color{codepurple},
    basicstyle=\ttfamily\footnotesize,
    breakatwhitespace=false,         
    breaklines=true,                 
    captionpos=b,                    
    keepspaces=true,                 
    numbers=left,                    
    numbersep=5pt,                  
    showspaces=false,                
    showstringspaces=false,
    showtabs=false,                  
    tabsize=2
}

\lstset{style=mystyle}

\title{Software Documentation}
\author{Your Name}
\date{\today}

\begin{document}

\maketitle

\tableofcontents
\newpage

\section{Introduction}
This document provides detailed documentation for the software package, describing classes used for modeling power systems components.

\section{Classes}

\subsection{Bus Class}

\subsubsection{Overview}
The \texttt{Bus} class models an electrical bus in a power network.

\subsubsection{Class Definition}
\begin{lstlisting}[language=Python, caption=Bus Class Code]
class Bus:
    counter = 0
    
    def __init__(self, name: str, baseKV: float, vpu: float = 1, delta: float = 0, bus_type: str = 'Slack'):
        ...
        
    def __repr__(self):
        ...
\end{lstlisting}

\subsubsection{Attributes}
\begin{itemize}
    \item \texttt{name} (str): Name of the bus.
    \item \texttt{baseKV} (float): Base voltage level (kV).
    \item \texttt{vpu} (float): Per-unit voltage magnitude.
    \item \texttt{delta} (float): Voltage angle (degrees).
    \item \texttt{real\_power} (float): Real power injection.
    \item \texttt{reactive\_power} (float): Reactive power injection.
    \item \texttt{type} (str): Bus type ("Slack", "PQ", "PV").
    \item \texttt{index} (int): Unique identifier for the bus.
\end{itemize}

\subsubsection{Methods}
\begin{itemize}
    \item \texttt{\_\_init\_\_()}: Constructor, validates input and assigns properties.
    \item \texttt{\_\_repr\_\_()}: Returns readable representation.
\end{itemize}

\subsection{Conductor Class}

\subsubsection{Overview}
The \texttt{Conductor} class models an electrical conductor's physical and electrical properties.

\subsubsection{Class Definition}
\begin{lstlisting}[language=Python, caption=Conductor Class Code]
class Conductor:
    def __init__(self, name: str, diameter: float, GMR: float, resistance: float, ampacity: float):
        ...
\end{lstlisting}

\subsubsection{Attributes}
\begin{itemize}
    \item \texttt{name} (str): Name of the conductor.
    \item \texttt{diameter} (float): Diameter (feet or specified unit).
    \item \texttt{GMR} (float): Geometric Mean Radius.
    \item \texttt{resistance} (float): Electrical resistance (ohms).
    \item \texttt{ampacity} (float): Maximum current-carrying capacity (amps).
\end{itemize}

\subsubsection{Methods}
\begin{itemize}
    \item \texttt{\_\_init\_\_()}: Constructor.
\end{itemize}

\subsection{Bundle Class}

\subsubsection{Overview}
The \texttt{Bundle} class models a bundle of multiple conductors for high-voltage transmission.

\subsubsection{Class Definition}
\begin{lstlisting}[language=Python, caption=Bundle Class Code]
class Bundle:
    def __init__(self, name: str, num_conductors: int, spacing: float, conductor: Conductor):
        ...
        
    def calculate_dsc(self):
        ...
        
    def calculate_dsl(self):
        ...
        
    def __repr__(self):
        ...
\end{lstlisting}

\subsubsection{Attributes}
\begin{itemize}
    \item \texttt{name} (str): Name of the bundle.
    \item \texttt{num\_conductors} (int): Number of conductors in the bundle.
    \item \texttt{spacing} (float): Spacing between conductors (feet).
    \item \texttt{conductor} (\texttt{Conductor}): Associated conductor object.
\end{itemize}

\subsubsection{Methods}
\begin{itemize}
    \item \texttt{calculate\_dsc()}: Calculates the bundle's equivalent radius (\( D_{sc} \)).
    \item \texttt{calculate\_dsl()}: Calculates the bundle's self-GMR (\( D_{sl} \)).
    \item \texttt{\_\_repr\_\_()}: Provides a readable summary.
\end{itemize}

\subsection{Generator Class}

\subsubsection{Overview}
The \texttt{Generator} class models a generator attached to a bus in a power network.

\subsubsection{Class Definition}
\begin{lstlisting}[language=Python, caption=Generator Class Code]
class Generator:
    def __init__(self, name: str, bus, voltage_setpoint: float, mw_setpoint: float):
        ...
\end{lstlisting}

\subsubsection{Attributes}
\begin{itemize}
    \item \texttt{name} (str): Name of the generator.
    \item \texttt{bus} (Bus): Bus object to which the generator is connected.
    \item \texttt{voltage\_setpoint} (float): Desired bus voltage (p.u.).
    \item \texttt{mw\_setpoint} (float): Real power output (MW).
\end{itemize}

\subsubsection{Methods}
\begin{itemize}
    \item \texttt{\_\_init\_\_()}: Constructor.
\end{itemize}

\section{Examples}

\subsection{Example: Creating a Bus and Generator}
\begin{lstlisting}[language=Python]
b1 = Bus(name="Bus 1", baseKV=138.0, vpu=1.0, delta=0.0, bus_type='Slack')
g1 = Generator(name="Gen 1", bus=b1, voltage_setpoint=1.0, mw_setpoint=100.0)
print(b1)
print(g1)
\end{lstlisting}

\subsection{Example: Creating a Bundle}
\begin{lstlisting}[language=Python]
c1 = Conductor(name="Drake", diameter=1.108, GMR=0.460, resistance=0.028, ampacity=845)
bundle = Bundle(name="Bundle 1", num_conductors=2, spacing=18, conductor=c1)
print(bundle)
\end{lstlisting}

\end{document}